\section{Discussion}\label{sec:discussion}
Our retrospective evaluations expose a discrepancy between strong performance on 09B and substantially weaker performance on Future-B. The hybrid retains a higher fraction of top-1 accuracy than several comparator models across the two cohorts, yet even its later accuracy remains below 50\%. Robustness is therefore relative, not immunity to temporal degradation.

The difference between micro-only and hybrid performance is compatible with the hypothesis that a complementary representation helps classification. However, an apparent hybrid advantage may also depend on the fusion mechanism, an additional classifier, feature engineering, or training choices. We will isolate these factors through the existing frozen ablation evidence and, where necessary, further Historical-DEV-only experiments. Causal language about longitudinal information should be reserved for contrasts that actually hold the relevant alternatives fixed.

A second consideration concerns deployment. If macro features depend on labeled historical observations of each monitored website, the attacker must be capable of maintaining such a reference history. If the macro branch uses information from a queried visit's future or from labels unavailable at inference time, its operational interpretation would be substantially narrower. We will audit the construction window, labeling assumptions, and update rules in the code before making any deployment claim.

The CUMUL-style and Tik-Tok baselines illustrate that matching the chronological protocol does not automatically make architectures and preprocessing choices equivalent to their original publications. In particular, the CUMUL-style byte representation and the WFlib-adapted NetCLR settings require transparent deviation tables. Our goal is a fair, inspectable benchmark with appropriately qualified scope rather than an unqualified claim of state-of-the-art superiority.
